\section{Further research questions and proposed program}
\label{sec:questions}

The following questions are ordered roughly from immediate engineering tasks to
longer-term topology.  They are phrased so that a negative answer would also be
informative.

\subsection{Immediate repository work}

\begin{enumerate}[label=\textbf{Q\arabic*.}]
\item \textbf{Find a genuine accumulated-support separator.}
      Does there exist a PD-code unknot for which the reviewed last-touch search
      fails at a given depth but the one-birth accumulated-footprint search
      succeeds at that depth?  An exhaustive search over small diagrams would
      turn Proposition~\ref{prop:last-vs-accumulated} from an abstract warning
      into a knot-specific regression fixture.

\item \textbf{Find a genuine two-birth RIII unlock.}
      Is there an unknot diagram whose shortest first-unlocking RIII trace has
      birth number at least two with the conservative dart footprint?  If so,
      minimize crossings and depth and add the result as an adversarial test.

\item \textbf{Measure the empirical parameter distribution.}
      On the existing random closures, scrambled unknots, hard survivors, and
      named examples, record the minimum successful pair $(b,k)$ under a fixed
      trial cap.  Separate successful and failed search costs, and publish raw
      traces rather than only aggregate counts.

\item \textbf{Incremental candidate maintenance.}
      Can legal RIII triangles meeting the active footprint be updated in
      amortized constant time per move, rather than rediscovered from nearby
      darts?  The dart involution changes only a bounded footprint, so a local
      invalidation table should be possible.

\item \textbf{Canonical commutation reduction.}
      Beyond sorting births, can independent connected moves be quotiented by a
      trace-monoid normal form?  A safe partial-order reduction could remove many
      permutations without changing completeness.
\end{enumerate}

\subsection{Parameterized topology}

\begin{enumerate}[label=\textbf{Q\arabic*.},resume]
\item \textbf{Bounded births for natural diagram classes.}
      Are reduced alternating unknot diagrams (necessarily trivial after
      reduction), positive braid-derived unknots, arborescent unknots, or
      diagrams of bounded projection-treewidth guaranteed to have bounded birth
      number during monotone simplification?

\item \textbf{Logarithmic depth from bounded width.}
      Does bounded treewidth, branchwidth, carving width, or pathwidth of the
      projection graph imply $k=O(\log N)$ after an appropriate choice of
      crossing-decreasing reduction order?  A balanced-separator proof would
      directly combine with the present theorem.

\item \textbf{Birth number versus dependency-tree width.}
      The support-overlap graph of a shortest trace is connected.  Relate the
      minimum birth number to standard directed or undirected parameters of this
      graph: number of sources under chronological orientation, pathwidth,
      vertex separation, or search number.  Which parameter admits the best
      deterministic enumeration?

\item \textbf{Lower bounds.}
      Construct families of unknot diagrams requiring
      $\Omega(\log N)$, $\Omega(N^\epsilon)$, or $\Omega(N)$ births or RIII
      depth under crossing-nonincreasing RI/RII/RIII simplification.  Such a
      family would sharply delimit this route.

\item \textbf{Reduction-order optimization.}
      The current simplifier takes available RI/RII moves greedily.  Can a
      different choice of decreasing move reduce the later birth/depth
      parameters?  Determine the complexity of finding an optimal reduction
      order and whether local confluence or matroid-like exchange principles
      apply.
\end{enumerate}

\subsection{Algorithmic improvements to the birth exponent}

\begin{enumerate}[label=\textbf{Q\arabic*.},resume]
\item \textbf{Color-coding independent births.}
      Can randomized or derandomized color-coding find $b$ disjoint causal birth
      footprints in $2^{O(b)}\poly(N)$ rather than $N^b$ time, while still
      accounting for their later joins?

\item \textbf{Terminal-rooted reverse search.}
      RIII is an involution.  Can one enumerate causal predecessors backward
      from potential RI/RII footprints, meeting a forward search in the middle?
      The obstacle is that the final diagram is unknown; a bounded local-state
      interface might overcome it.

\item \textbf{Representative families.}
      When several birth sets induce the same boundary behavior on an active
      footprint, can all but a small representative family be discarded?  This
      would parallel rank-based or connectivity-based dynamic programming.

\item \textbf{Planar localization.}
      Prove that a shortest birth footprint can be chosen inside a bounded-radius
      neighborhood of the eventual connecting action, or find a counterexample.
      A positive theorem would replace global $N$ choices by a function of $k$.

\item \textbf{Learning only as a scheduler.}
      Can a learned model predict which birth candidates to try first while the
      complete bounded search remains available underneath?  The learned layer
      must affect order only, never pruning or verdicts.
\end{enumerate}

\subsection{Combination with exact invariants}

\begin{enumerate}[label=\textbf{Q\arabic*.},resume]
\item \textbf{Khovanov-aware RIII ordering.}
      Can inexpensive frontier estimates predict which legal RIII move will
      reduce later Khovanov scan width, not merely crossing count?  The synthesis
      already observes that fewer crossings can increase scan work.

\item \textbf{Joint parameter theorem.}
      Suppose RIII birth/depth are moderately large but the resulting diagram
      has a small exact-scan frontier.  Derive a combined bound that branches on
      RIII only until a certified frontier threshold is reached, then switches
      to exact contraction.

\item \textbf{Twist-run compression after causal simplification.}
      Does a small RIII causal footprint preserve or reveal a low signed-braid
      run count?  If so, the twist-compression theorem and the present theorem
      could be composed rather than run as unrelated alternatives.

\item \textbf{Hierarchy transfer.}
      Relate a bounded RIII causal region to a bounded modification of a handle
      hierarchy or boundary pattern.  This could provide a three-dimensional
      certificate that replaces a large remaining diagram search.
\end{enumerate}

\subsection{Formal verification}

\begin{enumerate}[label=\textbf{Q\arabic*.},resume]
\item \textbf{Lean formalization of the abstract theorem.}
      Formalize strong disjoint commutation, causal connectedness, birth
      front-loading, and the counting bound.  The Boolean prototype supplies a
      finite executable semantics, but the proof should be independent of it.

\item \textbf{Verified dart footprint.}
      Prove that the concrete \code{apply\_r3} function writes only the claimed
      24-dart footprint and is an involution on every validated diagram.

\item \textbf{Proof-carrying traces.}
      Specify a stable JSON schema for RIII/RI/RII traces, source hashes, budgets,
      and replay results.  Make the replay kernel the only trusted component for
      positive simplification verdicts.
\end{enumerate}

\subsection{A focused next milestone}

The most informative next experiment is not to raise depth globally.  It is to
implement opt-in $b=1$ accumulated support, preserve the current budget, and run
paired tests on the exact 117-depth corpus recorded by the synthesis.  Three
outcomes are all useful:
\begin{itemize}
\item new successes establish that last-touch locality was a real completeness
      bottleneck;
\item no new successes but low overhead suggests the stronger search can safely
      replace the heuristic;
\item high failed-search overhead identifies the need for partial-order
      reduction or a better scheduler before increasing $b$.
\end{itemize}
